% Source: https://artofproblemsolving.com/community/c3h3180799_2017_emcc_individual_speed_test__exeter_math_club_competition
% Stable identity: aops-topic-3180799-q19
Place $B=(0,0),A=(20,0),C=(0,37),D=(17,37)$. Consider the midpoint $Q=(\dfrac{37}{2},\dfrac{37}{2})$ of $AD$. A perpendicular to $AD$ has slope $\dfrac{3}{37}$, so through this point it meets $BC$ at $P=(0,17)$. The quadrilateral $B A Q P$ has area $\dfrac{1}{2}(20\cdot\dfrac{37}{2}+17\cdot\dfrac{37}{2})=\dfrac{1369}{4}$, exactly half of the trapezoid area $\dfrac{1}{2}(20+17)37=\dfrac{1369}{2}$. Parallel cuts in this direction increase the area on one side strictly as they move across the trapezoid, so this is the unique equal-area cut. Therefore $\boxed{BP=17}$.
\begin{asy}
usepackage("amsmath");
size(7cm);
pair A=(20,0),B=(0,0),C=(0,37),D=(17,37),P=(0,17),Q=(18.5,18.5);draw(A--B--C--D--cycle);draw(P--Q);label("$A$",A,SE);label("$B$",B,SW);label("$C$",C,NW);label("$D$",D,NE);label("$P$",P,W);label("$Q$",Q,E);
\end{asy}
